% ---------------------------------------------------------------------
%  Rövidítések és jelölések
% ---------------------------------------------------------------------
\chapter*{Rövidítések és jelölések}
\addcontentsline{toc}{chapter}{Rövidítések és jelölések}
\thispagestyle{plain}

\begin{singlespace}
\subsection*{Rövidítések}
\begin{longtable}{@{}p{3.4cm}p{11.4cm}@{}}
AQG & \emph{Automatic Question Generation}, automatikus kérdésgenerálás\\
BM25 & \emph{Best Matching 25}, valószínűségi relevanciakeretre épülő lexikális rangsoroló függvény\\
CI & \emph{Confidence Interval}, konfidenciaintervallum\\
E5 & \emph{EmbEddings from bidirEctional Encoder rEpresentations}, kontrasztív szövegbeágyazó modellcsalád\\
GGUF & a \texttt{llama.cpp} kvantált modellformátuma\\
GQA & \emph{Grouped-Query Attention}, csoportosított lekérdezésű figyelem\\
GPU & \emph{Graphics Processing Unit}, grafikus feldolgozóegység\\
GDPR & az EU 2016/679 általános adatvédelmi rendelete\\
KV & \emph{key--value} (kulcs--érték) gyorsítótár a transzformer dekódolásakor\\
LLM & \emph{Large Language Model}, nagy nyelvi modell\\
MCQ & \emph{Multiple-Choice Question}, feleletválasztós kérdés\\
MRR & \emph{Mean Reciprocal Rank}, átlagos reciprok rang\\
Q4\_K\_M & 4 bites, csoportos K-kvantálás közepes változata (\texttt{llama.cpp})\\
RAG & \emph{Retrieval-Augmented Generation}, visszakereséssel kiegészített generálás\\
RQ & \emph{Research Question}, kutatási kérdés\\
RRF & \emph{Reciprocal Rank Fusion}, reciprok rangfúzió\\
SLM & \emph{Small Language Model}, kis nyelvi modell (itt: $\le$\,8 milliárd paraméter)\\
VRAM & a GPU dedikált memóriája\\
\end{longtable}

\subsection*{Fontosabb jelölések}
\begin{longtable}{@{}p{3.4cm}p{11.4cm}@{}}
$\Doc$ & a forrásdokumentum (karaktersorozat)\\
$S=(s_1,\dots,s_m)$ & a dokumentum mondatai\\
$\Chunks=(c_1,\dots,c_K)$ & a dokumentum darabjai (\emph{chunk}), $K=|\Chunks|$\\
$\phi(\cdot)\in\R^{d}$ & egységnyi normájú szövegbeágyazás ($d=768$, multilingual-E5-base)\\
$\delta_i$ & a szomszédos mondatok koszinusz-távolsága, $\delta_i = 1-\langle \phi(s_i),\phi(s_{i+1})\rangle$\\
$\tau$ & a darabolási alapküszöb (percentilis vagy $\mu+k\sigma$)\\
$\sigma=(n,T,\lambda,\ell)$ & vizsgaspecifikáció: kérdésszám, típusok, nehézség, nyelv\\
$q=(t,O,o^\star,R)$ & kérdés: törzs, válaszlehetőségek, kulcs, hivatkozott darabok\\
$\Exam=(q_1,\dots,q_n)$ & a generált vizsga\\
$\mathcal{B}=\{(\kappa_j,T_j,b_j)\}$ & a vizsgaterv slotjai: fogalom, típus, Bloom-szint\\
$G(q),\,A(q),\,V(q)$ & alátámasztottság, vak megoldhatóság, disztraktor-validitás (bírói ítéletek)\\
$F(\Exam)$ & formai megfelelés indikátora\\
$\mathrm{Cov}(\Exam)$ & darablefedettség\\
$Q(\Exam)$ & minőségi index, a négy ráta átlaga\\
$L(q)$ & válaszszivárgás indikátora (kulcs a törzsben)\\
$\kappa$ & Cohen-féle egyezési együttható (a szövegkörnyezetből egyértelmű, hogy nem fogalom)\\
$\alpha$ & Krippendorff-féle megbízhatósági együttható\\
$r_{rb}$ & páros rangbiszeriális hatásméret\\
\end{longtable}
\end{singlespace}
